\section{Hypothesis development}
\label{sec:hypotheses}

The empirical analysis tests eight pre-specified main hypotheses
(H1--H14, restricted to the eight with non-trivial findings) and
seven post-specification extensions (C1--C7). This section motivates
the eight main hypotheses from theory and prior evidence; the full
list with formal test specifications appears in
Section~\ref{ssec:hypothesis-tests}.

\subsection{Magnitude response (H1)}

Under semi-strong market efficiency \citep{fama1970efficient}, the
arrival of an unscheduled informational event should produce a
price revision that is detectable in the conditional second moment
of returns even when its conditional mean is small or sign-ambiguous.
\citet{andersen2003micro} demonstrate this directly for scheduled
macro announcements in high-frequency FX, and
\citet{calomiris2019news} document analogous volatility effects for
unscheduled news in equity markets cross-country. We hypothesise
that, on intraday windows from one to two hundred and forty
minutes, the absolute event-window return exceeds the absolute
matched-baseline return after controlling for hour-of-day and
day-of-week seasonality. This is H1.

\subsection{Directional content of sentiment labels (H2)}

The textual-analysis tradition since \citet{tetlock2007giving}
documents that sentiment extracted from news text predicts future
returns, though the directional signal is small in absolute terms.
\citet{heston2017news} report that neural sentiment classifiers
extract richer directional content than dictionary classifiers
\citep{loughran2011liability}, particularly for negative news. We
hypothesise that the discrete bull/bear LLM sentiment predicts the
realised target direction at hit rates above 50\%, with effect sizes
in the single-digit percentage points consistent with prior
literature. This is H2.

\subsection{Sentiment-trend interaction (H3)}

The classical reversal literature \citep{tetlock2007giving} suggests
that the directional content of sentiment depends on prior price
behaviour: bullish news following a price drop produces a larger
upward revision than bullish news following a price run-up, because
the latter has been at least partially priced in. We hypothesise
that the interaction between LLM sentiment and the prior 60-minute
trend is a significant predictor of the absolute event-window
return, controlling for the sentiment and trend main effects. This
is H3.

\subsection{Closed-period gap (H4)}

When markets are closed (weekends, holidays), news that arrives
during the closure should be impounded into the opening price
discontinuously rather than continuously. The cumulative target
sentiment across the closure should therefore predict the realised
opening-bar gap, in line with \citet{fama1970efficient} and standard
event-study results on overnight returns. This is H4.

\subsection{Pre-event drift and information arrival timing (H8)}

The textual-analysis literature typically assumes that the public
publication time of a news headline coincides with the
informational event time, which is appropriate for
machine-readable wire feeds with millisecond timestamps. For
real-time chat-aggregated public feeds, the latency between the
primary source and the public publication is variable and can be
substantial. We hypothesise that the absolute return in the
fifteen-minute window preceding the public timestamp exceeds the
matched baseline at a magnitude comparable to the post-event drift,
consistent with primary-source latency rather than with insider
trading or systematic front-running. This is H8.

\subsection{Persistence (H9)}

If the news effect represents a permanent re-pricing rather than a
transient liquidity shock, the sign of the short-horizon return
should agree with the sign of the longer-horizon return at a rate
above 50\%. The classical asset-pricing literature
\citep{fama1970efficient,mackinlay1997event} does not distinguish
these mechanisms at intraday frequencies. We hypothesise that the
sign of the $+15$-minute target return agrees with the sign of the
$+4$-hour target return at a rate above $50\%$, with the magnitude
of the longer-horizon return larger than the magnitude of the
shorter-horizon return. This is H9.

\subsection{Time-of-day heterogeneity (H11)}

Intraday volatility is strongly heterogeneous across the trading
day, with established peaks around major exchange openings and
macro-release windows \citep{andersen2003micro}. The event/baseline
magnitude ratio should reflect this heterogeneity if our
matched-baseline construction succeeds at absorbing the
intraday cycle. We hypothesise that hour-of-day effects on the
event/baseline ratio are statistically detectable for at least one
of the two target assets. This is H11.

\subsection{Cross-asset relation (H14)}

Macro-news events typically affect USD and equity markets through
correlated risk-on/risk-off mechanisms. \citet{andersen2003micro}
document positive correlation between USD strength and S\&P 500
returns conditional on macro news. We hypothesise that, under the
USD-proxy convention, the event-window USD return and the
event-window NDX return are positively correlated. This is H14.

\subsection{Summary}

The eight main hypotheses cover the four dimensions on which an
event-study of unscheduled online news should provide evidence:
magnitude (H1), direction (H2 and H3), timing (H4, H8, H9), and
cross-sectional structure (H11, H14). The remaining tests in the
pipeline (H5--H7, H10, H12--H13 and C1--C7) are auxiliary or
robustness extensions, motivated below where they are reported.
